Eight well known historical supernovae will appear in the projected sky one
at a time (not necessarily in chronological order). You have to identify the
appropriate map (Map 1 / Map 2) where a particular supernova belongs and
mark it in the corresponding map with a `+' sign and write codes `S1' to
`S8' beside it.

Each supernova code will be projected on the dome for 10 seconds, followed
by appearance of the supernova for 60 seconds and then 20 seconds for you to
mark the answers.

\begin{description}
\item[(OP1.1)] For S1, S2, S3, S4 and S5, the projected sky corresponds to
  the sky as seen from Rio de Janeiro on the midnight of 21st May.
\item[(OP1.2)] For S6, S7 and S8, the projected sky corresponds to the sky
  as seen from Beijing on the midnight of 20th November. There will be a gap
  of two minutes after S5 for change over and adaptation to the new sky.
\end{description}

% FIGURE NEEDED: the two answer sky maps (Map 1 and Map 2) on which the
% supernova positions are to be marked
